\documentclass[10pt,a4paper]{amsart}
\usepackage[margin=19mm]{geometry}
\usepackage{amsmath,amssymb,mathtools}
\usepackage[scr=rsfs,cal=euler]{mathalfa}
\usepackage{xfrac}
\usepackage[unicode,hidelinks,bookmarks=false]{hyperref}
\newcommand{\ie}{i.e.\ }
\newcommand{\cf}{cf.\ }
\newcommand{\resp}{respectively\ }
\newcommand{\Subsection}{\subsection}
\providecommand{\nobreakdash}{\nobreak\hbox{-}}
\setlength{\emergencystretch}{2em}
\multlinegap0pt
\usepackage{bm}
\usepackage{bbm}
\usepackage{dsfont}
\usepackage{xspace}
\usepackage{cancel}
\usepackage{mathtools}
\usepackage[shortlabels]{enumitem}
\usepackage{amsthm}
\makeatletter
\@ifundefined{theorem}{\newtheorem{theorem}{Theorem}}{}
\@ifundefined{lemma}{\newtheorem{lemma}[theorem]{Lemma}}{}
\makeatother
\newcommand{\C}{\mathbb{C}}
\newcommand{\N}{\mathbb{N}}
\newcommand{\R}{\mathbb{R}}
\newcommand{\supp}{\mbox{supp}\,}
\title[From zero solutions with partially bounded supports to solva]{From zero solutions with partially bounded supports to solvability and Runge approximation for partial differential equations}
\author{Tomasz Ci\'as, Thomas Kalmes}
\date{Source: arXiv:2607.19272v1 (CC BY 4.0)}
\begin{document}
\maketitle

\noindent\emph{Source and licence.} From zero solutions with partially bounded supports to solvability and Runge approximation for partial differential equations. Version arXiv:2607.19272v1 (CC BY 4.0), distributed under the Creative Commons Attribution 4.0 International licence, as recorded in the arXiv licence metadata for this exact version and shown on the abstract page https://arxiv.org/abs/2607.19272v1. Full rights notice: licenses/arxiv-2607.19272-CC-BY-4.0.txt.\par\smallskip

\noindent\textit{Editorial scope.} Counted unit: the Theorem whose source label is \texttt{th:zero solution standard form} in the article (source0.tex lines 309--319, source label \texttt{th:zero solution standard form}), delivered together with the article's own complete original proof of it (source0.tex lines 321--408, one \texttt{proof} environment of 88 lines). No mathematics is written, repaired or re-derived: statement and proof are the article's own text line for line. Required context retained uncounted: cohoon-solution (lines 173--184); lemma:convolution f ast g (lines 283--289). Every definition, display and label of those lines is retained unchanged. Omitted from this package and not redistributed: 1--172, 185--282, 290--308, 320--320, 409--1851. Nothing inside a retained excerpt is omitted or altered: every retained body is delivered exactly as the article's lines are written, so no omission receipt of any kind accompanies this package. Citations: the retained excerpts carry no citation command at all, so no bibliography entry of the article is transcribed and no reference is redistributed. Reference adaptation: none was needed and none was made; every reference inside the delivered unit resolves inside it, so no unresolved reference or citation is produced. Printed slips kept literal: no printed word, symbol, hypothesis or number of the counted statement or of its proof was altered. Layout and class: the article's own class is \texttt{amsart} with packages including amsmath, amssymb, latexsym, amscd, fontenc, lmodern, mathrsfs, amsthm, amsfonts, url, mathtools, enumerate, graphicx, color; the wrapper is the lane's pinned \texttt{amsart} editorial wrapper at 10pt on A4, which restates the theorem-like environments the retained text needs and adds the article's own macro definitions that the retained text needs (\texttt{\textbackslash C}, \texttt{\textbackslash N}, \texttt{\textbackslash R}, \texttt{\textbackslash supp}), with extra wrapper packages (bm, bbm, dsfont, xspace, cancel, mathtools, enumitem). The delivered numbering is the unit's own. Assets: no retained span contains a figure environment, a graphics-inclusion command, a PDF-page inclusion, a table or a TikZ picture; no image file is copied into this package, the package declares no asset dependency and no image file is redistributed with it. Licence: version arXiv:2607.19272v1 of this article is distributed under the Creative Commons Attribution 4.0 International licence, as recorded in the arXiv licence metadata for this exact version and shown on the abstract page https://arxiv.org/abs/2607.19272v1. A full rights notice is delivered at licenses/arxiv-2607.19272-CC-BY-4.0.txt. Characters: every retained excerpt is pure ASCII, so the delivered engine, XeLaTeX with its default Latin Modern text font, reports no missing character.
\begin{theorem}\label{cohoon-solution}
	Let $m\in\N$ and let $P\in\C[\xi_1,\ldots,\xi_d]$ be a polynomial of the form $P(\xi',\xi_d)=\xi_d^m+\sum_{k=0}^{m-1}Q_k(\xi') \xi_d^k$, where $Q_k\in\C[\xi_1,\ldots,\xi_{d-1}]$ with $\operatorname{deg}(Q_k)\leq m-1-k$, $k=0,\ldots, m-1$. Moreover, let $Y\subset\R^{d-1}$ be open and let $s\in[1,\infty)^{d-1}$ be such that $\max_k s_k<\frac{m}{m-1}$. For every $h_0, \ldots, h_{m-1}\in G^s(Y)$ there is $u\in G^{(s_1,\ldots,s_{d-1},1)}(Y\times\R)$ such that the following properties hold:
	\begin{enumerate}[{\upshape(i)}]
		\item $u\in\mathscr{E}_P(Y\times\R)$, 
		\item $\partial^{j}_{d}u(x',0)=h_j(x')$ for all $x'\in Y$, $0\leq j\leq m-1$,
		\item $\supp u\subset\left(\cup_{j=0}^{m-1}\supp h_j\right)\times\R$.
	\end{enumerate}
	Especially, for every $x'\in Y$ the function $u(x',\cdot)$ is real analytic on $\R$, and $u(\cdot, x_d)\in G^s(Y)$ for every $x_d\in\R$.
	Consequently , whenever $\sigma>1$, for all $\epsilon>0$ there is $u\in G^{(\sigma,\ldots,\sigma,1)}(\R^{d-1}\times \R)$ satisfying $P(D)u=0$ such that
	\[ \supp u= \overline{B_{d-1}(0,\epsilon)}\times\R\] 
    and $u$ is real analytic on $B_{d-1}(0,\epsilon)\times\R$, where for $n\in\N$ we denote the euclidean ball in $\R^n$ of radius $\rho$ centered at the origin by $B_n(0,\rho)$.
\end{theorem}

\begin{lemma}\label{lemma:convolution f ast g}
	Let $f,g\in \mathscr{E}(\R^d)$ with $\supp f\subset\prod_{k=1}^d I_k$ and $\supp g\subset\prod_{k=1}^d J_k$ for some open (not necessarily bounded) intervals $I_k,J_k\subset\R$. Let us assume that, for all $1\leq k\leq d$, $I_k$ or $J_k$ is bounded. Then, the following assertions hold.
	\begin{enumerate}[{\upshape(i)}]
		\item The convolution $f\ast g$ is well-defined.
		\item $f\ast g\in \mathscr{E}(\R^d)$ and $D^{\alpha}(f\ast g)=(D^{\alpha}f)\ast g=f\ast (D^{\alpha}g)$ for all $\alpha\in\N_0^d$.
	\end{enumerate}
\end{lemma}

\begin{theorem}\label{th:zero solution standard form}
	Let $m\in\N, 1\leq l\leq d-1$, and let $P\in\C[\xi_1,\ldots,\xi_d]$ be a polynomial of degree $m$ with principal part $P_m$ of the form 
    \[P_m(\xi)=\xi_d^m+\sum_{j=l+1}^{d-1}p_j(\xi)\xi_j\]
    for some polynomials $p_j\in\C[\xi_1,\ldots,\xi_d]$.
    Then, for every $0<\delta<\epsilon$ there is $v\in\mathscr{E}_P(\R^d)$ such that
	\begin {enumerate}[{\upshape(i)}]
	\item $\overline{B_{l}(0,\delta)}\times\R^{d-l}\subset\operatorname{supp} v\subset \overline{B_{l}(0,\epsilon)}\times\R^{d-l}$,
	\item $v$ is real analytic on $B_{l}(0,\delta)\times\R^{d-l}$.
	\end{enumerate} 
    As before, for $n\in\N$ we denote the euclidean ball in $\R^n$ of radius $\rho>0$ centered at the origin by $B_n(0,\rho)$.
\end{theorem}

\begin{proof}
Let $P(\xi)=P_m(\xi)+\sum_{|\alpha|<m}c_\alpha\xi^\alpha$. Set $\epsilon_1=\frac{\epsilon-\delta}{2}$ and $\epsilon_2=\frac{\epsilon+\delta}{2}$. Next, let us fix a nonnegative function $\lambda\in \mathscr{E}(\R^{d-1})$ such that $\lambda$ depends only on the first $l$ coordinates,
\[\supp\lambda=\overline{B_l(0,\epsilon_2)}\times\R^{d-l-1},\] and such that $\lambda$ is even symmetric, and $\lambda$ is real analytic on $B_l(0,\epsilon_2)\times\R^{d-l-1}$. 
	
	Let $1<s<\frac{m}{m-1}$ and let us fix a nonzero and nonnegative function $\phi\in G^s(\R^{d-1})$ with $\supp\phi=\overline{B_{d-1}(0,\epsilon_1)}$ and $\phi>0$ in $B_{d-1}(0,\epsilon_1)$.
	By Theorem \ref{cohoon-solution} applied for $h_0=\phi, h_1=\ldots,h_{m-1}=0$, there is a solution $u\in \mathscr{E}(\R^d)$ of the partial differential equation
	\[\left(D_d^m+\sum_{|\alpha|<m} c_\alpha D^\alpha\right)u=0\]
	such that (see also the proof of the last part of Theorem \ref{cohoon-solution})
	$$\supp{u}=\supp\phi\times\R=\overline{B_{d-1}(0,\epsilon_1)}\times\R\subset[-\epsilon_1,\epsilon_1]^{d-1}\times\R,$$
	$u(x',0)=\phi(x')$ for all $x'\in\R^{d-1}$, and $u(x',\cdot)$ is real analytic on $\R$. In particular, it can be extended to a function $u\colon\R^{d-1}\times\C\to\C$ in such a way that $u(x',\cdot)$ is an entire function.
	
	By a standard application of Lebesgue's Dominated Convergence Theorem to differentiability of parameter intergrals, the function
	\[f\colon\C\to\C, f(z):=\int_{\left[-\epsilon_1,\epsilon_1\right]^{d-1}}u(x',z)\lambda(x')\mathrm{d}x',\]
	is entire and since
	\begin{align*}
		f(0)=\int_{\left[-\epsilon_1,\epsilon_1\right]^{d-1}}u(x',0)\lambda(x')\mathrm{d}x'=\int_{\left[-\epsilon_1,\epsilon_1\right]^{d-1}}\phi(x')\lambda(x')\mathrm{d}x'>0,
	\end{align*}
	$f$ is nonzero. 
	
	Let $\mu\in \mathscr{E}(\R)$ be a nonzero, nonnegative function such that $\supp\mu\subset\left[-\epsilon_2,\epsilon_2\right]$ and $\mu$ is real analytic on $\left(-\epsilon_2,\epsilon_2\right)$. We set $g\colon\R\to\C$, $g(t):=\overline{f(-t)}\mu(t)$. Then $g\in \mathscr{E}(\R)$, $\supp g\subset\left[-\epsilon_2,\epsilon_2\right]$ and $g$ is real analytic on $\left(-\epsilon_2,\epsilon_2\right)$. With these, let us also define
	\[\psi\colon\R^{d-1}\times\R\to\R, \psi(x',x_d):=\lambda(x')g(x_d).\]
	Since $\lambda$ depends only on the first $l$ coordinates, $D_j\psi=0$ for $l+1\leq j\leq d-1$,
	\[\supp{\psi}\subset\overline{B_{l}(0,\epsilon_2)}\times\R^{d-l-1}\times\left[-\epsilon_2,\epsilon_2\right]\] and $\psi$ is real analytic on
    \[B_{l}(0,\epsilon_2)\times\R^{d-l-1}\times\left(-\epsilon_2,\epsilon_2\right).\]
	
	Finally, we set $v:=u\ast\psi$. By Lemma \ref{lemma:convolution f ast g}, $v$ is well-defined and $v\in \mathscr{E}(\R^d)$. Moreover,
    \begin{align*}
		\supp{v}&\subset\supp u+\supp\psi\\
		&\subset\overline{B_{d-1}(0,\epsilon_1)}\times\R
		+\overline{B_{l}(0,\epsilon_2)}\times\R^{d-l-1}\times\left[-\epsilon_2,\epsilon_2\right]\\
		&\subset\overline{B_{l}(0,\epsilon)}\times\R^{d-l},
	\end{align*}
	and
	\begin{align*}
		v(0)&=\int_{\R^d}u(y)\psi(-y)\mathrm{d}y\\
		&=\int_{[-\epsilon_2,\epsilon_2]}\left(\int_{[-\epsilon_1,\epsilon_1]^{d-1}}u(y',y_d)\lambda(y')\mathrm{d}y'\right)g(-y_d)\mathrm{d}y_d\\
		&=\int_{[-\epsilon_2,\epsilon_2]}f(y_d)g(-y_d)\mathrm{d}y_d=\int_{[-\epsilon_2,\epsilon_2]}f(y_d)\overline{f(y_d)}\mu(y_d)\mathrm{d}y_d\\
		&=\int_{[-\epsilon_2,\epsilon_2]}|f(y_d)|^2\mu(y_d)\mathrm{d}y_d>0
	\end{align*}
	as well as
	\begin{align*}
		P(D)v&=\left(D_d^m+\sum_{j=l+1}^{d-1}p_j(D)D_j+\sum_{|\alpha|<m}c_\alpha D^\alpha\right)v\\
		&=\left(D_d^m+\sum_{|\alpha|<m}c_\alpha D^\alpha\right)v
		+\sum_{j=l+1}^{d-1}p_j(D)D_jv\\
		&=\left[\left(D_d^m+\sum_{|\alpha|<m}c_\alpha D^\alpha\right)u\right]\ast\psi+u\ast\sum_{j=l+1}^{d-1}p_j(D)D_j\psi=0.
	\end{align*}
	
	It remains to show that $v$ is real analytic on $B_l(0,\delta)\times\R^{d-l}$ which by $v(0)>0$ and the identity principle for real analytic functions will also entail
	$$\overline{B_{l}(0,\delta)}\times\R^{d-l}\subset \supp v. $$
	Let $K$ be a compact subset of $B_l(0,\delta)\times\R^{d-l}$. Denoting by $\pi_d\colon \R^d\rightarrow\R$ the projection onto the $d$-th axis, let $R>0$ be such that $\pi_d(K)\subset(-R,R)$. Because
	\[\supp u= \overline{B_{d-1}(0,\epsilon_1)}\times\R\]
	and
	\[\supp\psi(x-\cdot)\subset x-\overline{B_l(0,\epsilon_2)}\times\R^{d-l-1}\times\left[-\epsilon_2,\epsilon_2\right],\]
	for $x\in K$ and $\alpha,\beta\in\N_0^d$ it holds
    \begin{align*}
        \supp \partial^\alpha u\cap\supp\partial^\beta\psi(x-\cdot)& \subset\overline{B_{d-1}(0,\epsilon_2)}\times [-R-\epsilon_2,R+\epsilon_2]\\
        &\subset [-\epsilon_1,\epsilon_1]^{d-1}\times [-R-\epsilon_2,R+\epsilon_2].
    \end{align*}
	Hence, for $x\in K$ and $\alpha\in\N_0^d$,
	\begin{eqnarray}\label{eq:real analyticity 1}
		|\partial^\alpha v(x)|
		&=&\left|\left(\partial^{(0,\ldots,0,\alpha_d)}u\ast\partial^{(\alpha_1,\ldots,\alpha_{d-1},0)}\psi\right)(x)\right|\nonumber\\
		&\leq&\int_{[-\epsilon_1,\epsilon_1]^{d-1}\times[-R-\epsilon_2,R+\epsilon_2]}|\partial^{(0,\ldots,0,\alpha_d)}u(y)|\cdot\left|\partial^{(\alpha_1,\ldots,\alpha_{d-1},0)}\psi(x-y)\right|\mathrm{d}y\nonumber\\
		&\leq& (2\epsilon_1)^{d-1}2(R+\epsilon_2) \max_{y\in [-\epsilon_1,\epsilon_1]^{d-1}\times[-R-\epsilon_2,R+\epsilon_2]}\left|\partial^{(0,\ldots,0,\alpha_d)}u(y)\right|\times\nonumber\\
		&&\times \max_{y\in [-\epsilon_1,\epsilon_1]^{d-1}\times[-R-\epsilon_2,R+\epsilon_2]}\left|\partial^{(\alpha_1,\ldots,\alpha_{d-1},0)}\psi(x-y)\right|.
	\end{eqnarray}
	By the Cauchy integral formula, for each $y=(y',y_d)\in \overline{B_{d-1}(0,\epsilon_1)}\times[-R-\epsilon_2,R+\epsilon_2]$ we have
	\begin{eqnarray}\label{eq:real analyticity 2}
		\left|\partial^{(0,\ldots,0,\alpha_d)}u(y)\right|
		&=&\frac{1}{2\pi}\alpha_d!\left|\int_{|\xi|=R+\epsilon}\frac{u(y',\xi)}{(\xi-y_d)^{\alpha_d+1}}\mathrm{d}\xi\right|\nonumber\\
		&\leq& \left(\max\left\{\frac{1}{\epsilon_1},1\right\}\right)^{\alpha_d+1}(R+\epsilon)\\
        &\times&\max\left\{|u(y',\xi)|\colon y'\in\left[-\epsilon_1,\epsilon_1\right]^{d-1},|\xi|=R+\epsilon\right\}\alpha_d!.\nonumber
	\end{eqnarray}
	
	Because for $x\in K$ and $y\in \overline{B_{d-1}(0,\epsilon_1)}\times[-R-\epsilon_2,R+\epsilon_2]$ it holds
    \[x-y\in B_{l}(0,\delta+\epsilon_1)\times \R^{d-l}\]
	and since $\lambda$ is real analytic in $B_{l}(0,\epsilon_2)\times\R^{d-l-1}$, there is $C_3>0$ such that for every $x=(x',x_d)\in K$ and $y=(y',y_d)\in \overline{B_{d-1}(0,\epsilon_1)}\times[-R-\epsilon_2,R+\epsilon_2]$ we have
	$$|\partial^{\alpha'}\lambda(x'-y')|\leq C_3^{|\alpha'|}\alpha'!$$
	for each $\alpha'\in\N_0^{d-1}$. Consequently,
	\begin{eqnarray}\label{eq:real analyticity 3}
		&&\max_{x\in K, y\in \overline{B_{d-1}(0,\epsilon_1)}\times[-R-\epsilon_2,R+\epsilon_2]}\left|\partial^{(\alpha_1,\ldots,\alpha_{d-1},0)}\psi(x-y)\right|\nonumber\\
		&=&\max_{x\in K, y\in \overline{B_{d-1}(0,\epsilon_1)}\times[-R-\epsilon_2,R+\epsilon_2]}\left|\partial^{(\alpha_1,\ldots,\alpha_{d-1})}\lambda(x'-y')\right|\left|g(x_d-y_d)\right|\nonumber\\
		&\leq&C_3^{\alpha_1+\cdots\alpha_{d-1}}\alpha_1!\cdots\alpha_{d-1}!\max_{t\in [-\epsilon_2,\epsilon_2]} \left|g(t)\right|.
	\end{eqnarray}
	Combining \eqref{eq:real analyticity 1}, \eqref{eq:real analyticity 2}, and \eqref{eq:real analyticity 3}, we obtain $M, C>0$ such that for every $x\in K$ and $\alpha\in\N_0^d$
	\[|\partial^\alpha v(x)|\leq C^{|\alpha|}\alpha! M\]
	which finally shows that $v$ is real analytic on $B_{l}(0,\delta)\times\R^{d-l}$. 
\end{proof}

\end{document}
